\documentclass[11pt,a4paper]{article}
\usepackage[utf8]{inputenc}
\usepackage[T1]{fontenc}
\usepackage[spanish]{babel}
\usepackage{hyperref}
\usepackage{geometry}
\geometry{margin=2.5cm}
\usepackage{parskip}

\title{El estandar}
\author{Alejandro Valenzuela Rivera\\\small Universidad Autónoma de Nuevo León}
\date{\today}
\begin{document}
\maketitle

\begin{quote}
\textbf{Sesgo declarado.} Este trabajo se sostiene desde el sesgo declarado de su autor, y no finge neutralidad. Lo que afirma se puede tumbar, y ese es el punto: el falsificador está escrito, no escondido.
\end{quote}

\section{La tesis}
\textbf{\textquestiondown{}Se puede escribir un estándar que no envejezca con la era?}

\section{Por qué se sostiene}
\textbf{Porque su contexto refutable está escrito, y se puede tumbar sin que esta pregunta caiga.} \textbf{Y porque el estándar entero es la prueba: seis capacidades, seis verbos, y ninguno nombra lenguaje ni dominio.}

\section{De dónde viene}
\textbf{De `Tesis/ILUMINAR.md`, y de las correcciones de la autoridad: «todo el estándar está mal, el estándar es atemporal porque metes python; esta mierda va a envejecer mal». Y después: «los entes también son atemporales; todo faro es atemporal; mantis es de la era».}

\section{Lo que queda abierto}
\textbf{\textquestiondown{}Qué distingue un estándar de una moda con reglas?}

\section{El caso}
\textbf{La prueba no es una obra: es el `RUTA.pr` de una consola de mano y el de un brazo.} \textbf{Si la capacidad se llama igual en los dos, el estándar no envejeció.}

\section{Método}
\textbf{\textquestiondown{}Cada capacidad del mundo nombra un acto, y no una herramienta? \textquestiondown{}Cada verbo del ejecutor se cumple sin que nadie sepa qué lenguaje lo corre? \textquestiondown{}Se puede escribir una tarea que no diga `python` por ningún lado?}

\section{Límites, y qué la tumba}
\textbf{Un estándar que nombra su lenguaje; una capacidad que dice `git` en vez de «historial»; un verbo que solo se puede cumplir con una herramienta; un estándar que cambia cuando cambia la industria.}

\section{Estado de la tesis}
\textbf{escrita, y medida.} \textbf{`piso.py` corre seis capacidades y ninguna nombra lenguaje. `juez.py` corre seis verbos y ningún caso sabe qué los corre.} ---

\section*{Cómo se leyó esta tesis}
Esto se escribió en el repositorio man-made-empathy, y la tesis está completa en vision del autor/tesis/man made empathy/el estandar/. La máquina que la compuso no agregó nada que no estuviera escrito: si algo de esto es falso, se dice en el mismo sitio, y se corrige con fecha.
\end{document}
